\scanpage{94}
A)
\[
\left\{\begin{aligned}
G&=\mathrm{Sl}(2,\Lambda)=\left\{\begin{pmatrix}a&b\\c&d\end{pmatrix}\in M_2(\Lambda)\mid ad-bc=1\right\}\\
H&=\left\{\begin{pmatrix}1&\deleted{\unsure{\lambda}}t\\0&1\end{pmatrix}\mid\deleted{\lambda}t\in\Lambda\right\}
\end{aligned}\right.\tag*{(49)}
\]

Soit $g$
\[
\begin{aligned}
g&=\begin{pmatrix}a&b\\c&d\end{pmatrix}\in G,\quad\text{d'où}\\
g^{-1}&=\begin{pmatrix}d&-b\\-c&a\end{pmatrix}\qquad\text{on a}
\end{aligned}
\]
\[
\operatorname{int}(g)\begin{pmatrix}1&\lambda\\0&1\end{pmatrix}=
\underbrace{\begin{pmatrix}a&b\\c&d\end{pmatrix}\begin{pmatrix}1&\deleted{\lambda}t\\0&1\end{pmatrix}}_{\begin{pmatrix}a&a\deleted{\lambda}t+b\\c&c\deleted{\lambda}t+d\end{pmatrix}}
\begin{pmatrix}d&-b\\-c&a\end{pmatrix}=
\begin{pmatrix}1\deleted{\unsure{+}}-ac\deleted{\lambda}t&a^2\deleted{\lambda}t\\-c^2\deleted{\lambda}t&1+ac\deleted{\lambda}t\end{pmatrix}\tag*{(*)}
\]
donc
\[
\operatorname{int}(g)\begin{pmatrix}1&\deleted{\lambda}t\\0&1\end{pmatrix}\in H
\Longleftrightarrow\left\{\begin{aligned}
-c^2\deleted{\lambda}t&=0\quad\text{d'où }(1-ac\lambda)(1+ac\lambda)=1\\
\deleted{\unsure{ac\lambda=0}}\\[-3pt]
1-ac\deleted{\lambda}t&=1
\end{aligned}\right.
\]
Donc
\[
N_G(H)=\left\{g=\begin{pmatrix}a&b\\c&d\end{pmatrix}\in\mathrm{Sl}(2,\Lambda)\ \middle|\ 
\begin{aligned}-c^2\deleted{\lambda}t&=0\\
\deleted{\unsure{1-}}ac\deleted{\lambda}t&=0\end{aligned}\right\}
\deleted{\unsure{\forall\lambda\in\Lambda\}}}
\]
Or \deleted{faisant $\lambda=1$, cela signifie $c^2=0$, \inserted{\unsure{et $ac$}} donc}
comme $ad-bc=1$ i.e. $ad=1+bc$, on voit
que \inserted{$c^2=0\Longrightarrow$} $ad$ est inversible \deleted{($(ad)^{-1}=1-bc$ avec} avec
\[
(ad)^{\unsure{-1|+1}}=1-bc\qquad a^{-1}=(1-bc)d
\]
et comme $ac=0$, on doit avoir $c=0$

Donc
\[
\begin{aligned}
N_G(H)&=\left\{g=\begin{pmatrix}\deleted{\unsure{a}}\lambda&b\\0&\lambda^{-1}\end{pmatrix}\mid\lambda\in\Lambda^*,\ b\in\Lambda\right\}\\
&\shortparallel\\[-3pt]
&\left\{g\in\mathrm{Sl}(2,\Lambda)\mid\operatorname{int}(g)\begin{pmatrix}1&1\\0&1\end{pmatrix}\in H\right\}
\end{aligned}\tag*{(50)}
\]
Donc une \deleted{\unsure{autom}} permutation de $M^*$ qui
commute à $\mathrm{Sl}(2,\Lambda)$ transforme $x$ en $\lambda x$, \inserted{$(\lambda\in\Lambda^*)$}
donc c'est l'homothétie $\lambda\,\mathrm{id}_M$, qui a le
$=$ effet sur $x$ et commute à $\mathrm{Sl}(2,\Lambda)$.
% archive sheet number: 93
